% Abstract types: representation independence. Verified: Q1 stores the queue
% in order and appends; Q2 stores it reversed and conses; a client sees no
% difference.
\part[3]\hfill

  Both structures below seal to the same signature, but represent a queue
  differently: \code{Q1} keeps its elements front-to-back, while \code{Q2} keeps
  them \textbf{reversed}.
  \begin{codeblock}
    signature QUEUE =
      sig
        type 'a t
        val empty  : 'a t
        val push   : 'a -> 'a t -> 'a t
        val toList : 'a t -> 'a list     (* front to back *)
      end

    structure Q1 :> QUEUE = struct
      type 'a t = 'a list
      val empty = []
      fun push x q = q @ [x]
      fun toList q = q
    end

    structure Q2 :> QUEUE = struct
      type 'a t = 'a list
      val empty = []
      fun push x q = x :: q
      fun toList q = rev q
    end
  \end{codeblock}

  Can a client written against \code{QUEUE} tell the two apart? Explain.

  \begin{solutionorbox}[16em]
    No. A client can only build queues with \code{empty} and \code{push} and can
    only observe them with \code{toList}, and both structures agree on every such
    sequence of operations --- e.g.\ pushing \code{1}, \code{2}, \code{3} onto
    \code{empty} and calling \code{toList} gives \code{[1, 2, 3]} in both.

    The differing internals are invisible because \code{'a t} is
    \textbf{abstract}: the client cannot inspect a queue directly, only through
    the operations the signature provides. This is what lets an implementation be
    swapped for another --- say, a faster one --- without any client noticing.
  \end{solutionorbox}
